\def\title{Determining k Clusters in k-Means Clustering with the CRAB Algorithm}
\def\author{Jasmine Kristine Soriano Cabrera}
\def\degreemonth{June}
\def\degreeyear{2026}
\def\degree{Master of Science in Statistics}
\def\numberofmembers{3}
   \def\chair{Kelly Bodwin, Ph.D. \linebreak Professor of Statistics}
   \def\othermemberA{Allison Theobold, Ph.D. \linebreak Professor of Statistics}
   \def\othermemberB{Julia Schedler, Ph.D. \linebreak Professor of Statistics}
\def\keywords{Clustering, Data Science, k-Means.} % Make sure to end this in a . (i.e. Some, Thing, Interesting.)

\def\abstract{%
Unsupervised clustering algorithms such as k-Means and Hierarchical Clustering
are commonly used in research and academia. While many different scoring methods
have been developed to determine the most optimal number of clusters (i.e.
the Elbow Method, Silhouette Score, etc.), their vague results and initial
assumptions about the data, present the need for a new approach in finding the
optimal k number of clusters. To address this, a
new cluster scoring method, the CRAB algorithm, was created. This approach works
specifically with k-Means and utilizes random subsampling of the data along with
the classification of its pairwise points. The constructed definition of “Buddies”
and “Rivals” provides insight into which observation points tend to stick together
and which tend to be further apart. For research and educational purposes, an R
package is being built, implementing different functions to help understand the
many capabilities of the cRab algorithm.

}

\def\tocpagelegnth{1} % Number of pages your Table of Contents is (used to set correct page numbers for List of Tables and List of Figures)
\def\totflag{1} % 1 if you have at least 1 table, otherwise 0 (adds Table of Tables page)
\def\tofflag{1} % 1 if you have at least 1 figure, otherwise 0 (adds Table of Figures page)
